\section{Ellipsoid optimization and upper-bound certificates}\label{sec:2}

The same optimization problem appears twice: first when choosing a center, and again when recomputing the largest inscribed ellipsoid after a cut. We begin with its variables and constraints, then derive the short inequality used throughout the paper.

\subsection{From containment to a convex problem}\label{sec:2.1}

Throughout, the dimension is a positive integer. Let $B_2^n=\{z:\|z\|_2\le1\}$ be the closed Euclidean unit ball. Write
\begin{equation}\label{eq:2.1}
\mathcal E(c,A)=c+AB_2^n,\qquad A\in\mathbb S_{++}^n,
\qquad \operatorname{vol}\mathcal E(c,A)=\kappa_n\det A,
\quad \kappa_n=\operatorname{vol}(B_2^n).
\end{equation}

The matrix $A$ is the linear map applied to the ball. Its eigenvalues are the semiaxis lengths. For a single inequality $g^Tx\le\beta$, containment is equivalent to
\begin{equation}\label{eq:2.2}
\max_{\|z\|\le1}g^T(c+Az)=g^Tc+\|Ag\|\le\beta.
\end{equation}

Thus, for a bounded full-dimensional polytope
$\mathcal P=\{x:g_r^Tx\le\beta_r,\ r=1,\ldots,N\}$, the ellipsoid problem is
\begin{equation}\label{eq:2.3}
\begin{aligned}
\text{maximize}_{c,\ A\succ0}\quad&\log\det A\\
\text{subject to}\quad&g_r^Tc+\|Ag_r\|\le\beta_r,
\qquad r=1,\ldots,N.
\end{aligned}
\end{equation}

The objective is concave and the constraints are convex \cite[Section 8.4.2]{BV2004}. Each support constraint is also a second-order-cone constraint:
$(\beta_r-g_r^Tc,Ag_r)\in\mathcal L_{n+1}$, where
$\mathcal L_{n+1}=\{(\tau,z):\tau\ge\|z\|\}$.
We remove redundant inequalities with $g_r=0$.

\subsection{A certificate cancels the unknown center}\label{sec:2.2}

To bound the volume of every feasible ellipsoid, assign a scalar $\alpha_r$ and
a vector $q_r$ to each supporting inequality. With
$\operatorname{sym}X=(X+X^T)/2$, the complete certificate conditions are
\begin{equation}\label{eq:2.4}
\begin{gathered}
\sum_r\alpha_rg_r=0,\qquad \sum_r\alpha_r\beta_r=n,\\
\|q_r\|\le\alpha_r\quad\text{for every }r,\\
M=\sum_r\operatorname{sym}(q_rg_r^T)\succ0.
\end{gathered}
\end{equation}

Balance cancels the unknown center, and the second equality fixes the scale.
The norm inequality is the second-order-cone constraint
$(\alpha_r,q_r)\in\mathcal L_{n+1}$; it also implies $\alpha_r\ge0$.
The final condition collects the directional information into a positive-definite
matrix. Thus \eqref{eq:2.4} includes every constraint used by the dual problem below.

\begin{lemma}[{volume certificate}]\label{lem-certificate}
If \eqref{eq:2.4} holds, every feasible ellipsoid in \eqref{eq:2.3} satisfies
\begin{equation}\label{eq:2.5}
\operatorname{tr}(AM)\le n,\qquad
\det A\le\frac1{\det M}.
\end{equation}

\end{lemma}

\begin{proof}
The two conclusions require two steps: first turn the support
constraints into a trace bound, then apply scalar AM--GM to obtain a determinant
bound. Aggregate the support inequalities; balance cancels the center:
\begin{equation}\label{eq:2.6}
\begin{aligned}
\operatorname{tr}(AM)
&=\sum_r q_r^TAg_r
\le\sum_r\alpha_r\|Ag_r\|\\
&\le\sum_r\alpha_r(\beta_r-g_r^Tc)
=n-c^T\sum_r\alpha_rg_r=n.
\end{aligned}
\end{equation}

To pass from trace to determinant, use the symmetric positive-definite matrix
$Y=M^{1/2}AM^{1/2}$. Its trace and determinant are the quantities we need.
For its eigenvalues $\rho_1,\ldots,\rho_n>0$, scalar AM--GM gives
\begin{equation}\label{eq:2.7}
\det A\det M=\det Y=\prod_{j=1}^n\rho_j
\le\left(\frac1n\sum_{j=1}^n\rho_j\right)^n
=\left(\frac{\operatorname{tr}(AM)}n\right)^n\le1.
\end{equation}

Dividing by $\det M>0$ proves \eqref{eq:2.5}. 
\end{proof}

The matrix $M$ is a certificate: checking its finite constraints is enough to bound the optimum. A larger determinant gives a stronger bound. Optimizing the certificate leads to the normalized dual problem
\begin{equation}\label{eq:2.8}
\text{maximize }\log\det M
\quad\text{over }(\alpha_r,q_r)\text{ satisfying \eqref{eq:2.4}}.
\end{equation}

\subsection{Exact certificates and optimal contacts}\label{sec:2.3}

Optimizing the certificate recovers the ellipsoid optimum exactly. We prove
this by constructing multipliers from the contacts of the optimal ellipsoid.
The only geometric ingredient is the following classical characterization.

\begin{lemma}[{John's contact characterization \cite{John1948} and \cite[Theorems 3.1 and 3.3]{GG2007}}]\label{lem-john}
Let $K\subset\mathbb R^n$ be a convex body containing $B_2^n$.
Then $J(K)=B_2^n$ if and only if there are finitely many unit vectors
$u_i\in\partial K\cap S^{n-1}$ and positive weights $\lambda_i$ such that
\begin{equation}\label{eq:2.9}
u_i^Tx\le1\ (x\in K),\qquad
\sum_i\lambda_i u_i=0,\qquad
\sum_i\lambda_i u_i u_i^T=I,
\qquad \sum_i\lambda_i=n.
\end{equation}

\end{lemma}

We use this standard result without proof. The last identity follows by taking
traces in the second-moment identity. We call finite unit vectors and positive
weights with the two moment identities a \emph{John frame}; the inequalities in
\eqref{eq:2.9} say that the frame supports $K$. The same vector $u_i$ specifies both
the point of tangency with the unit ball and the outward normal of its
supporting plane.

\textbf{Construct an optimal certificate.} Return to the bounded full-dimensional
polytope $\mathcal P$ of \eqref{eq:2.3}. The next proposition proves attainment and
equality of the primal and normalized dual values.

\begin{proposition}[{optimal multipliers from John contacts}]\label{prop-exact-duality}
Problems \eqref{eq:2.3} and \eqref{eq:2.8} attain their optima. If $\mathcal E(c,A)=J(\mathcal P)$, there are normalized dual multipliers with aggregate $M=A^{-1}$ and zero cone pairings. Hence the optimal values of \eqref{eq:2.3} and \eqref{eq:2.8} sum to zero.

\end{proposition}

\begin{proof}
Construct the dual variables from the contacts of the actual
optimal ellipsoid. We first identify the active facet normal for each contact,
then rescale the John weights so that their moments become the dual constraints.

\textbf{Identify the active facets.} Let $\mathcal E(c,A)=J(\mathcal P)$,
whose existence is classical, and apply John's characterization to
$A^{-1}(\mathcal P-c)$. A unit contact $p$ lies on some defining plane.
For its index $r$,
\begin{equation}\label{eq:2.10}
\|Ag_r\|\le\beta_r-g_r^Tc=(Ag_r)^Tp
\le\|Ag_r\|\|p\|=\|Ag_r\|.
\end{equation}

Hence $p=Ag_r/\|Ag_r\|$ and the plane is active. Assign contacts to such indices and combine their weights. Writing $\ell_r=\|Ag_r\|>0$, this gives
\begin{equation}\label{eq:2.11}
p_r=\frac{Ag_r}{\ell_r},\qquad
\sum_r\mu_rp_r=0,\qquad
\sum_r\mu_rp_rp_r^T=I,\qquad \sum_r\mu_r=n.
\end{equation}

\textbf{Turn contact weights into multipliers.} The weights satisfy $\mu_r\ge0$ and
$\mu_r(\beta_r-g_r^Tc-\ell_r)=0$. Dividing by the active support length
$\ell_r$ makes the zero first moment become dual balance.
Define $\alpha_r=\mu_r/\ell_r$ and $q_r=\alpha_rp_r$. Then
\begin{equation}\label{eq:2.12}
\begin{aligned}
\sum_r\alpha_rg_r&=A^{-1}\sum_r\mu_rp_r=0,\\
\sum_r\alpha_r\beta_r
&=c^T\sum_r\alpha_rg_r+\sum_r\mu_r=n,\\
\|q_r\|&=\alpha_r,\\
\alpha_r(\beta_r-g_r^Tc)-q_r^TAg_r
&=\alpha_r(\beta_r-g_r^Tc-\ell_r)=0.
\end{aligned}
\end{equation}

\textbf{Use isotropy to match the inverse shape.} The second moment supplies the
remaining dual matrix equation. Since $g_r=\ell_rA^{-1}p_r$, the aggregate is
\begin{equation}\label{eq:2.13}
\sum_rq_rg_r^T=\sum_r\mu_rp_rp_r^TA^{-1}=A^{-1}.
\end{equation}

The cone constraints and zero pairings follow immediately. Lemma~\ref{lem-certificate} now proves that this normalized certificate is dual optimal and that the gap is zero. 
\end{proof}

\textbf{Recover contacts from an optimal pair.} The converse identifies the active
supports and shows that the normalized directions recover a John frame.

\begin{proposition}[{optimality and contact recovery}]\label{prop-kkt}
A feasible pair for \eqref{eq:2.3} and \eqref{eq:2.8} is optimal if and only if
\begin{equation}\label{eq:2.14}
M=A^{-1},\qquad
\alpha_r(\beta_r-g_r^Tc)-q_r^TAg_r=0\quad\text{for every }r.
\end{equation}

For each $\alpha_r>0$,
\begin{equation}\label{eq:2.15}
\beta_r-g_r^Tc=\|Ag_r\|,\qquad
q_r=\alpha_r\frac{Ag_r}{\|Ag_r\|}.
\end{equation}

With $p_r=Ag_r/\|Ag_r\|$ and $\mu_r=\alpha_r\|Ag_r\|$, the vectors and weights satisfy \eqref{eq:2.11}.

\end{proposition}

\begin{proof}
At a primal--dual optimum, Proposition~\ref{prop-exact-duality} gives
$\det A\det M=1$. Equality throughout the trace--AM--GM chain \eqref{eq:2.7} then forces
$Y=M^{1/2}AM^{1/2}=I$ and $\operatorname{tr}(AM)=n$.
The cone pairings are nonnegative and, by balance and normalization, their
sum is $n-\operatorname{tr}(AM)=0$. Thus every pairing vanishes.
Conversely, $M=A^{-1}$ gives equality in the certificate bound, proving
optimality of the feasible pair. For $\alpha_r>0$, equality in the support
and Cauchy--Schwarz inequalities gives \eqref{eq:2.15}.

The remaining point is to recover the full second moment from the symmetrized
aggregate. To recover isotropy, set $T=\sum_r\alpha_rg_rg_r^T/\|Ag_r\|$.
The aggregate equation is a Sylvester equation for $T$:
$AT+TA=2A^{-1}$. For $H=T-A^{-2}$, therefore,
\begin{equation}\label{eq:2.16}
AH+HA=0,
\qquad
0=\operatorname{tr}\bigl(H(AH+HA)\bigr)
=2\operatorname{tr}(AH^2)
\ge2\lambda_{\min}(A)\|H\|_F^2.
\end{equation}

Thus $H=0$ and $T=A^{-2}$. Positive definiteness has determined the
unsymmetrized moment matrix without a commutativity assumption.
The contact moments follow:
\begin{equation}\label{eq:2.17}
\sum_r\mu_rp_r=A\sum_r\alpha_rg_r=0,
\qquad
\sum_r\mu_rp_rp_r^T=ATA=I,
\qquad
\sum_r\mu_r=\operatorname{tr}I=n.
\end{equation}

This proves \eqref{eq:2.11}. 
\end{proof}

These results explain the exact optimization meaning of the certificate.
The universal contraction proof in Section~\ref{sec:3} uses Lemma~\ref{lem-certificate} and one explicit
feasible certificate; it does not require solving the normalized dual.
Section~\ref{sec:10.1} gives the usual unnormalized dual and its scale reduction.

\subsection{Finite contact reduction}\label{sec:2.4}

An invertible affine change of coordinates preserves ratios of ellipsoid
volumes. Send $J(K)$ to the unit ball and apply Lemma~\ref{lem-john}.
The resulting supporting frame converts the geometric problem into a finite
polytope problem while preserving the sharp supremum.

The frame does not need to describe the whole boundary of $K$. Retaining only
these supporting planes produces a polytope containing $K$, while John's
characterization certifies that its original maximal ellipsoid is unchanged.

\begin{lemma}[{finite contact reduction}]\label{lem-reduction}
Let $K$ be a convex body with
$J(K)=B_2^n$, and let $\mathcal U$ be a supporting John frame from Lemma~\ref{lem-john}.
Then the contact polytope defined below is a convex body satisfying
\begin{equation}\label{eq:2.18}
P_{\mathcal U}=\{x:u_i^Tx\le1\text{ for all }i\},\qquad
K\subseteq P_{\mathcal U},\qquad J(P_{\mathcal U})=B_2^n.
\end{equation}

In particular, $w(P_{\mathcal U})=w(K)=\kappa_n$. For every central halfspace
$H_v=\{x:v^Tx\le0\}$ with $\|v\|=1$, each ellipsoid in $K\cap H_v$ is also
feasible in $P_{\mathcal U}\cap H_v$, so
$w(K\cap H_v)\le w(P_{\mathcal U}\cap H_v)$.
Both cut sets have nonempty interior because they contain the corresponding
half of the unit ball.

\end{lemma}

\begin{proof}
The main point is boundedness: the supporting halfspaces could
otherwise define an unbounded set. We obtain a radius bound directly from
the two John moments. Since $\|u_i\|=1$, the polyhedron contains $B_2^n$.
For $x\in P_{\mathcal U}$, put $r=\|x\|$ and $s_i=u_i^Tx$. Then
$-r\le s_i\le1$, $\sum_i\lambda_i s_i=0$, and $\sum_i\lambda_i s_i^2=r^2$.
The product $(1-s_i)(r+s_i)$ is nonnegative and quadratic in $s_i$, so its
weighted sum can be evaluated using exactly these moments:
\begin{equation}\label{eq:2.19}
\begin{aligned}
0\le\sum_i\lambda_i(1-s_i)(r+s_i)
&=r\sum_i\lambda_i+(1-r)\sum_i\lambda_i s_i-\sum_i\lambda_i s_i^2\\
&=nr-r^2.
\end{aligned}
\end{equation}

Thus $r(r-n)\le0$ and $r\ge0$ imply $r\le n$. The finite intersection is closed and convex, so $B_2^n\subseteq P_{\mathcal U}\subseteq nB_2^n$ makes it a convex body. John's characterization gives $J(P_{\mathcal U})=B_2^n$, and the supporting inequalities give $K\subseteq P_{\mathcal U}$. 
\end{proof}

\textbf{Why the reduction preserves the sharp constant.} The enlargement can increase
the post-cut ellipsoid volume, but it leaves the pre-cut volume unchanged.
Consequently,
\begin{equation}\label{eq:2.20}
\frac{w(K\cap H_v)}{w(K)}
\le\frac{w(P_{\mathcal U}\cap H_v)}{w(K)}
=\frac{w(P_{\mathcal U}\cap H_v)}{w(P_{\mathcal U})}.
\end{equation}

For a particular $K$ the inequality can be strict. For the extremal problem,
however, the two classes have the same supremum. Every normalized convex body
admits such an enlargement. Conversely, any John frame defines a convex body
$P_{\mathcal U}$ with John ellipsoid $B_2^n$: the boundedness argument above
uses only the frame identities, and Lemma~\ref{lem-john} identifies its ellipsoid.
Thus contact polytopes themselves belong to the original class, and
\begin{equation}\label{eq:2.21}
C_n
=\sup_{\substack{\mathcal U\ \text{a John frame}\\\|v\|=1}}
\frac{w(P_{\mathcal U}\cap H_v)}{w(P_{\mathcal U})}
=\sup_{\substack{\mathcal U\ \text{a John frame}\\\|v\|=1}}
\frac{w(P_{\mathcal U}\cap H_v)}{\kappa_n}.
\end{equation}

The first direction gives an upper bound for every convex body; the converse
shows that no sharpness is lost. For each fixed frame and cut,
$w(P_{\mathcal U}\cap H_v)/\kappa_n$ is the maximum feasible determinant
in the finite ellipsoid problem \eqref{eq:2.3}, with supporting
inequalities $(g_i,\beta_i)=(u_i,1)$ and the one added cut $(g_0,\beta_0)=(v,0)$.
The geometric problem has therefore become a universal inequality for finite
contact data. It is this finite problem that the next section certifies.
